\color{unidarkgreen}{\subsection{Устройство резервирования отказа выключателя}}
\color{black}

\begin{enumerate}


%%%%%%%%%%%%%%%%%%%%%%%%%%%%%%%%%%%%%%%%%%%%%%%%%% Общие данные %%%%%%%%%%%%%%%%%%%%%%%%%%%%%%%%%%%%%%%%%%%%%%%%%%
\item Общие данные

\begin{enumerate}

\item
Функция УРОВ предназначена для быстрого и селективного отделения от энергосистемы поврежденного присоединения с отказавшим выключателем в аварийных ситуациях путем отключения смежных выключателей, и с возможностью предварительного повторного действия на отключение отказавшего выключателя.

\item
Логика работы устройства обеспечивает возможность реализации индивидуального УРОВ присоединения как в составе терминала ДЗШ (внутренний УРОВ), так и в составе терминалов защит и АУВ присоединений (внешний или распределенный УРОВ).
Выбор режима работы УРОВ осуществляется с помощью программного переключателя \textbf{«Режим УРОВ Вn»}.
При установленном положении \textbf{«Внеш.»} используется схема внешнего УРОВ, реализованная в терминалах защит и АУВ присоединений.
При положении \textbf{«Внут.»} используется схема внутреннего УРОВ, реализованная в терминале ДЗШ.

\end {enumerate}


%%%%%%%%%%%%%%%%%%%%%%%%%%%%%%%%%%%%%%%%%%%%%%%%%% Внутренний УРОВ %%%%%%%%%%%%%%%%%%%%%%%%%%%%%%%%%%%%%%%%%%%%%%%%%%
\item Внутренний УРОВ

\begin{enumerate}

\item
Внутренний УРОВ обеспечивает действие как на собственный выключатель («на себя»), так и на смежные выключатели.
Действие на смежные выключатели включает в себя:
\begin{itemize}
\item отключение смежных выключателей через защиты смежных присоединений;
\item останов ВЧ-передатчика защиты;
\item пуск аварийной команды ТО на противоположную сторону ЛЭП;
\item запрет АПВ выключателей.
\end{itemize}

\item
Пуск УРОВ осуществляется при срабатывании внутренних защит на отключение (\textit{«Отключение Bn»}), при наличии сигнала пуска УРОВ от внешних защит (\textit{«Пуск УРОВ Вn»}) или сигнала пуска УРОВ НН (\textit{«Пуск УРОВ НН Bn»}).

\item
Для подтверждения пуска УРОВ при КЗ предусмотрен контроль наличия тока через выключатель.
Контроль осуществляется с помощью ИО, ток срабатывания которого задается уставкой \textbf{«I УРОВ Вn»}.

\item
Для обеспечения надежного пуска УРОВ при кратковременном пропадании сигналов пуска до полного обесточивания отказавшего выключателя предусмотрен подхват от ИО УРОВ.

\item
Сигнал «Пуск УРОВ НН Вn» предусмотрен для пуска УРОВ от защит, срабатывание которых не сопровождается увеличением тока или сопровождается протеканием малых токов КЗ, недостаточных для срабатывания токовых ИО УРОВ.
Например, при пуске УРОВ от УРОВ НН при КЗ на стороне НН (авто)трансформатора за токоограничивающим реактором или при пуске УРОВ от газовой защиты при малой нагрузке или ее отсутствии.
Данный пуск осуществляется с контролем положения выключателя.
Фиксация включенного положения выключателя осуществляется по инверсному сигналу \textit{«РПО Вn»}.

\item
В устройстве предусмотрена возможность реализации двух принципов выполнения внутреннего УРОВ:
\begin{itemize}
\item с автоматической проверкой исправности выключателя;
\item с дублированным пуском.
\end{itemize}

\item
Для реализации принципа УРОВ с дублированным пуском необходимо задать следующие значения уставок: \textbf{«Контр. РПВ Вn = Введено»}, \textbf{«УРОВ Вn на себя = Выведено»}.
В этом случае УРОВ действует только на отключение смежных выключателей с дополнительным контролем положения «своего» выключателя.
Контроль осуществляется по сигналу от РПВ.
Действие УРОВ «на себя» может быть выполнено с контролем по тока при введенном программном ключе \textbf{«Контр. по I УРОВ Вn на себя»}.

\item
Для реализации принципа УРОВ с автоматической проверкой исправности выключателя необходимо задать следующие значения уставок: \textbf{«Контр. РПВ Вn = Выведено»}, \textbf{«УРОВ Вn на себя = Введено»}.
В это режиме УРОВ действует на отключение «своего» выключателя при наличии пусковых условий в течение выдержки времени \textbf{«Тср. УРОВ Вn на себя»}.
Повторное отключение позволяет исключить ложное срабатывание УРОВ на смежные выключатели за счет возврата ИО УРОВ после отключения выключателя.

\end {enumerate}


%%%%%%%%%%%%%%%%%%%%%%%%%%%%%%%%%%%%%%%%%%%%%%%%%% Внешний УРОВ %%%%%%%%%%%%%%%%%%%%%%%%%%%%%%%%%%%%%%%%%%%%%%%%%%
\item Внешний УРОВ

\begin{enumerate}

\item
Логика внешнего УРОВ предусматривает возможность приема сигналов срабатывания УРОВ от терминалов защит и АУВ присоединений.
Для этого используется внешний сигнал \textit{«Ср. УРОВ Вn»}.

\item
С целью повышения надёжности работы и исключения ложных срабатываний реализован токовый контроль, осуществляемый с помощью ИО внутренних УРОВ.

\item
Вывод функции УРОВ осуществляется либо программным переключателем \textbf{«Режим УРОВ Вn»}, либо посредством сигналов \textit{«Вывод Вn»} или \textit{«Вывод УРОВ Вn»}.

\item 
Уставки УРОВ выключателей приведены в таблицах \ref{app_set:urov1}-\ref{app_set:urov12}.

\item 
Алгоритм работы УРОВ В1 приведен на рисунке \ref{pic:UROV_Q1}, алгоритмы работы остальных выключателей выполнены аналогично.

\medskip
\begin{figure}[H]
\centering
\includegraphics[width=0.8\textwidth,height=0.8\textheight,keepaspectratio]{\fbpath/04.01 ДЗШ/661.3331 УРОВ/img/UROV_Q1.pdf}
\captionof{figure}{Функциональная схема логики УРОВ В1}\label{pic:UROV_Q1}
\end{figure}

\item
Сигналы срабатывания УРОВ СШ (\textit{«УРОВ 1С»} и \textit{«УРОВ 2С»}) формируются при срабатывании одного из внутренних УРОВ присоединений с учетом текущей фиксации присоединения с отказавшим выключателем или при наличии сигналов срабатывания внешнего централизованного УРОВ (\textit{«Центр. УРОВ 1С»} и \textit{«Центр. УРОВ 2С»}).
При нарушении фиксации присоединений сигналы срабатывания УРОВ формируются для обеих СШ.

\item
При длительном наличии сигналов пуска и срабатывания УРОВ более выдержки времени \textbf{«Тнеиспр. УРОВ»} формируется сигнал неисправности цепи УРОВ.

\item 
Алгоритм работы УРОВ СШ приведен на рисунке \ref{pic:UROV_SSH}. Уставки УРОВ СШ приведены в таблице \ref{app_set:urovsh}.

\medskip
\begin{figure}[H]
\centering
\includegraphics[width=0.8\textwidth,height=0.8\textheight,keepaspectratio]{\fbpath/04.01 ДЗШ/661.3331 УРОВ/img/UROV_SSH.pdf}
\captionof{figure}{Функциональная схема логики УРОВ СШ}\label{pic:UROV_SSH}
\end{figure}

\end {enumerate}

\end{enumerate}